\section{拉普拉斯变换的收敛域}\label{Convergence}

\subsection{收敛横坐标}\label{sub:converge_abscissa}
我们在定义拉普拉斯变换时，曾要求$|f(t)|\leq C \mathrm{e}^{at}$，这样拉普拉斯积分$\mathcal{L} f(s)=\int_{0}^{\infty}f(t)\mathrm{e}^{-st}\,dt$\\
在$\text{Re}\ s>a$上收敛，但是我们不知道这个收敛域是否是最精确的。本节就来讨论拉普拉斯变换收敛域的问题，并由此给出一些由拉普拉斯变换分析时域性质的方法。

\begin{definition}[收敛横坐标]
    \textbf{收敛横坐标}为使得拉普拉斯积分$$\mathcal{L} f(s)=\int_{0}^{\infty}f(t)\mathrm{e}^{-st}\,dt$$最大收敛域的收敛轴实部（横坐标）$\sigma_c$，使得$\mathcal{L} f(s)$恰好在$\text{Re}\ s>\sigma_c$上收敛，而在$\text{Re}\ s<\sigma_c$上发散，即
    \[\sigma_c=\inf\{a\in\mathbb{R}|\mathcal{L} f(s)\text{在}\text{Re}\ s>a\text{上收敛}\}\]
\end{definition}

为了寻找收敛横坐标，我们有两种思路：
\begin{itemize}
    \item 从时域出发，分析$|f(t)|$的增长速率；
    \item 从复频域出发，已知$F(s)$在$\text{Re}\ s>a$上是$f$的拉普拉斯变换，讨论它在更大的区域上是否也是$f$的拉普拉斯变换。
\end{itemize}

首先来看思路一。由于$f\in\mathcal{E} $，在有限区间上的表现不会影响拉普拉斯积分的收敛性，因此我们只需考虑$t\to\infty$时$f$的增长速率。仿照等价无穷大量，我们用$|f(t)|\sim C \mathrm{e}^{a t}$表示$|f(t)|$与$C \mathrm{e}^{a t}$在$t\to\infty$时具有相同增长速率，即
\[|f(t)|\sim C \mathrm{e}^{a t}\iff \lim_{t\to\infty}\frac{|f(t)|}{C \mathrm{e}^{a t}}=1\]
则\begin{align*}
    |f(t)|\sim C \mathrm{e}^{a t}\iff ln|f(t)|\sim ln C +a t\iff\frac{ln|f(t)|}{t}\sim a +\frac{ln C}{t}=a
\end{align*}

但是，我们很容易构造这样一个反例来说明这个方法并不总是有效。考虑$f(t)=\mathrm{e}^{t\sin t}$，则极限
$$\lim_{t\to\infty}\frac{ln|f(t)|}{t}=\lim_{t\to\infty}\sin t$$
不存在，但很明显$\mathcal{L} f(s)$在$\text{Re}\ s>1$上收敛。对照这个反例改进这一方法，得到：
\begin{theorem}
    设$f\in\mathcal{E} $，则收敛横坐标可由以下公式得到：
    \[\sigma_c=\varlimsup_{t\to\infty}\frac{ln|f(t)|}{t}\]
\end{theorem}
\begin{proof}
    \[\sigma_c=\varlimsup_{t\to\infty}\frac{ln|f(t)|}{t}\iff \forall \epsilon>0,\exists T>0,\text{s.t.}\forall t>T,\frac{ln|f(t)|}{t}<\sigma_c+\epsilon\]
    因此$t$充分大时，$|f(t)|<\mathrm{e}^{(\sigma_c+\epsilon)t}$，从而$\mathcal{L} f(s)$在$\text{Re}\ s>\sigma_c+\epsilon$上收敛，$\epsilon$任意，故$\mathcal{L} f(s)$在$\text{Re}\ s>\sigma_c$上收敛。

    另一方面，如果$\exists a' <\sigma_c$，使得$\mathcal{L} f(s)$在$\text{Re}\ s>a'$上收敛，则$\mathcal{L} f(s)$在$\frac{a'+\sigma_c}{2}$上收敛，因此
    \[\int_{0}^{\infty}|f(t)|\mathrm{e}^{-\frac{a'+\sigma_c}{2}t}\,dt<\infty\]
    由此可知，$|f(t)|\mathrm{e}^{-\frac{a'+\sigma_c}{2}t}\to 0,t\to\infty$，即$t$充分大时，$|f(t)|<\mathrm{e}^{\frac{a'+\sigma_c}{2}t}$，从而
    \[\varlimsup_{t\to\infty}\frac{ln|f(t)|}{t}\leq \frac{a'+\sigma_c}{2}<\sigma_c\]
\end{proof}

其实，以上的定理和证明对$\varlimsup_{t\to\infty}\dfrac{ln|f(t)|}{t}=\pm\infty$的情形也是成立的，它们分别对应$\mathcal{L} f(s)$无定义和$\mathcal{L} f(s)$为整函数的情形。

下面来看思路二。设$\mathcal{L} f(s)$在$\text{Re}\ s>a$上良定义且等于$F(s)$，我们讨论$F(s)$在更大的区域$\text{Re}\ s>a'(a'<a)$上是否也是$f$的拉普拉斯变换。我们在上一节的定理5.3.2.中处理过类似的情况，因此我们先来看$F(s)$满足这一定理的条件时的情形。

\begin{theorem}
    设$F(s)$是某个函数$f\in\mathcal{E} $的拉普拉斯变换，且
    $$|F(s)|\leq C|s|^{-\alpha},\alpha>1,\forall \text{Re}\ s>a'$$
    令
    \[\sigma_s=\sup\left\{\text{Re}\ s_j\in\mathbb{C}\left|s_j\text{为}F(s)\text{的奇点}\right.\right\}\]
    则有$\sigma_c=\sigma_s$。
\end{theorem}
\begin{proof}
    一方面，$\sigma_c\geq\sigma_s$。$\mathcal{L} f(s)$在$\text{Re}\ s<\sigma_s$上当然是发散的，因为$F(s)$在$\text{Re}\ s=\sigma_s$上有奇点。

    另一方面，$\sigma_c\leq\sigma_s$。用反证法。假设$\sigma_c>\sigma_s$，则在带状区域$\sigma_s<\text{Re}\ s\leq\sigma_c$上，$F(s)$满足定理5.3.2.的条件，它可以做拉普拉斯逆变换，得到唯一的时域函数，由于拉普拉斯逆变换
    $$\frac{1}{2\pi \mathrm{i}}\lim_{r\to\infty}\int_{b-ir}^{b+ir}F(s)\mathrm{e}^{st}\,ds$$
    不依赖于$b$的选取，该时域函数必然就是$f(t)$，从而$\mathcal{L} f(s)=F(s)$在$\text{Re}\ s>\sigma_s$上成立，矛盾。
\end{proof}

然而，很多函数不满足以上定理的条件，例如
$$f(t)=\mathrm{e}^{at},a\in\mathbb{R}^*,F(s)=\dfrac{1}{s-a}$$
（尽管对这类函数我们可以直接验证$\sigma_c=\sigma_s=a$）。下面给出一个类似的定理\cite{Laplace_Transform}，以说明我们的目标结论$\sigma_c=\sigma_s$的适用范围是相当大的，这样，我们就有很大的把握来通过$F(s)$的奇点直接分析$f(t)$的收敛横坐标。
\begin{theorem}
    设$F(s)$是某个函数$f\in\mathcal{E} $的拉普拉斯变换，且$f(t)$单调，仍然定义\[\sigma_s=\sup\left\{\text{Re}\ s_j\in\mathbb{C}\left|s_j\text{为}F(s)\text{的奇点}\right.\right\}\]
    则有$\sigma_c=\sigma_s$。
\end{theorem}
\begin{proof}
    一方面，我们已经说明$\sigma_c\geq\sigma_s$。

    另一方面，$\sigma_c\leq\sigma_s$。用反证法，假设$\sigma_c>\sigma_s$，任取$b>\sigma_c$，由于$F(s)$在$\text{Re}\ s>\sigma_s$上全纯，可以将它展开为$b$处的泰勒级数：
    \[F(s)=\sum_{n=0}^{\infty}\frac{F^{(n)}(b)}{n!}(s-b)^n,\text{Re}\ s>\sigma_s\]
    并且该级数的收敛半径至少为$b-\sigma_s$，因此在以下区域中，
    $$D(b,b-\sigma_s)=\{s\in\mathbb{C}||s-b|<b-\sigma_s\}$$
    该级数绝对收敛、内闭一致收敛。
    \begin{figure}[H]
        \centering
        \includegraphics[width=0.5\textwidth]{sigma_c}
        \caption{收敛域示意图}
    \end{figure}
    容易验证$f(t)(-t)^n \mathrm{e}^{-b t}$在$[0,\infty)$上的积分一致收敛，根据含参变量反常积分积分的求导法则，有
    \[F^{(n)}(b)=\int_{0}^{\infty}f(t)(-t)^n \mathrm{e}^{-b t}\,dt\]
    
    考虑截断积分
    \begin{align*}
        \int_{0}^{R}f(t)\mathrm{e}^{-a t}\,dt&=\int_{0}^{R}f(t)\mathrm{e}^{-b t}\sum_{n=0}^{\infty}\frac{(b-a)^n}{n!}t^n\,dt
    \end{align*}
    
    由$\sum_{n=0}^{\infty}\dfrac{(b-a)^n}{n!}t^n\,dt$在$[0,R]$上一致收敛，可交换求和与积分次序，得到
    \begin{align*}
        \int_{0}^{R}f(t)\mathrm{e}^{-a t}\,dt&=\sum_{n=0}^{\infty}\frac{(b-a)^n}{n!}\int_{0}^{R}f(t)t^n \mathrm{e}^{-b t}\,dt
    \end{align*}
    
    由于$f(t)$单调，不妨设$f(t)\geq 0$（可以考虑$g(t)=-f(t)$，$f(t)$单调递增，剩余的工作只是处理$u(t)$），而$b>a$，级数的每一项都是非负的。由单调收敛定理，有
    \begin{align*}
        \lim_{R\to\infty}\int_{0}^{R}f(t)\mathrm{e}^{-a t}\,dt&=\lim_{R\to\infty}\sum_{n=0}^{\infty}\frac{(b-a)^n}{n!}\int_{0}^{R}f(t)t^n \mathrm{e}^{-b t}\,dt\\
        &=\sum_{n=0}^{\infty}\frac{(b-a)^n}{n!}\int_{0}^{\infty}f(t)t^n \mathrm{e}^{-b t}\,dt\\
        &=\sum_{n=0}^{\infty}\frac{(b-a)^n}{n!}F^{(n)}(b)=F(a)
    \end{align*}
    即$\mathcal{L} f(a)=F(a)$收敛，从而在$\text{Re}\ s=a$上$\mathcal{L} f$也收敛：\begin{align*}
        |\mathcal{L} f(a+\mathrm{i}\omega)|&=\left|\int_{0}^{\infty}f(t)\mathrm{e}^{-(a+\mathrm{i}\omega) t}\,dt\right|\\
        &\leq \int_{0}^{\infty}\left|f(t)\mathrm{e}^{-a t}\right|\,dt\\
        &=\int_{0}^{\infty}f(t)\mathrm{e}^{-a t}\,dt&\left(f(t)\geq 0\right)\\
        &=\mathcal{L}f(a)<\infty
    \end{align*}
    由唯一延拓定理，$\mathcal{L} f(s)=F(s)$在$\text{Re}\ s>a$上成立。

    由$a$的任意性，我们证明了$\mathcal{L} f(s)=F(s)$在$\text{Re}\ s>\sigma_s$上成立，矛盾。
\end{proof}

为了说明定理的条件不能被去掉，我们给出以下反例\cite{Laplace_Transform}：
\begin{example}
    考虑函数
    \[f(t)=\frac{\sin(\mathrm{e}^t)}{\mathrm{e}^t}\]
    带入前文中的公式可得$\sigma_c=-1$。但是它的拉普拉斯变换为
    \begin{align*}
        F(s)&=\int_{0}^{\infty}\frac{\sin(\mathrm{e}^t)}{\mathrm{e}^t}\mathrm{e}^{-s t}\,dt\\
        &=\int_{1}^{\infty}\frac{\sin u}{u^{s+1}}\,du&(u=\mathrm{e}^t)\\
        &=-\evalat{\frac{\cos(u)}{u^{s+1}}}{1}{\infty}+\int_{1}^{\infty}\cos(u)\,d\frac{1}{u^{s+1}}\\
        &=\cos(1)-(s+1)\int_{1}^{\infty}\frac{\cos(u)}{u^{s+2}}\,du
    \end{align*}
    由A-D判别法可知，最后一行的积分在$\text{Re}\ s>-2$上收敛，因此$F(s)$可全纯延拓至$\text{Re}\ s>-2$，进一步，多次分部积分后，可以发现$F(s)$能够延拓为整函数。
\end{example}

总之，对于相当大范围的函数，我们可以\textbf{通过分析其拉普拉斯变换的奇点来确定收敛横坐标}$\sigma_c$，从而确定拉普拉斯变换的最大收敛域。在前面我们已经证明：$$f\in\mathcal{E}\implies \sigma_c=\varlimsup_{t\to\infty}\frac{ln|f(t)|}{t}$$
因此我们可以\textbf{通过分析拉普拉斯变换的奇点来确定$f(t)$的增长速率}，也可以\textbf{通过分析$f(t)$的增长速率来确定拉普拉斯变换的奇点位置}。

\subsection{系统的稳定性}\label{sub:stability_of_system}
基于以上结论，我们可以使用拉普拉斯变换分析系统的\textbf{稳定性}。在\ref{sub:concepts_of_system}中曾介绍，系统的稳定性是指：
\begin{definition}
    如果一个系统在激励信号有界时，响应也是有界的，则称该系统是稳定系统。
\end{definition}

对于线性时不变系统，稳定性有充要条件：
\begin{proposition}[稳定性的时域判则]
    设线性时不变系统的冲激响应为$h(t)$，则系统稳定的充分必要条件是$\int_{-\infty}^{\infty}|h(t)|\,dt<\infty$。
\end{proposition}
\begin{proof}
    一方面，对于有界输入$|e(t)|<M$，有
    \begin{align*}
        |r(t)|&=\left|\int_{-\infty}^{\infty}h(\tau)e(t-\tau)\,d\tau\right|\\
        &\leq \int_{-\infty}^{\infty}|h(\tau)||e(t-\tau)|\,d\tau\\
        &<M\int_{-\infty}^{\infty}|h(\tau)|\,d\tau<\infty
    \end{align*}
    因此输出有界。

    另一方面，设$h(t)$不绝对可积，取有界激励信号$e(t)=\text{sgn}(h(-t))$，则
    \begin{align*}
        |r(0)|&=\left|\int_{-\infty}^{\infty}h(\tau)e(-\tau)\,d\tau\right|\\
        &=\left|\int_{-\infty}^{\infty}h(\tau)\text{sgn}(h(\tau))\,d\tau\right|\\
        &=\int_{-\infty}^{\infty}|h(\tau)|\,d\tau=\infty
    \end{align*}
    输出无界。
\end{proof}

下面使用拉普拉斯变换来分析系统的稳定性。对因果的线性时不变系统，我们定义
\begin{definition}[系统函数]
    设$h\in\mathcal{E} $是线性时不变系统的冲激响应，则系统函数定义为
    $$H(s)=\frac{R(s)}{E(s)}=\mathcal{L} h(s)$$
    如果虚轴$\text{Re}\ s=0$在$H(s)$的收敛域内，则$H(\mathrm{i}\omega)$与第四章中定义的频率响应一致。
\end{definition}
对$r(t)=h(t)*e(t)$两边同时进行拉普拉斯变换，有
\[\mathcal{L} r(s)=\mathcal{L} h(s)\mathcal{L} e(s),\text{即}R(s)=H(s)E(s)\]

\begin{proposition}[稳定性的复频域判则]
    设因果线性时不变系统的冲激响应为$h(t)\in\mathcal{E}$，系统函数为$H(s)=\mathcal{L} h(s)$，则稳定性的充分必要条件是$H(s)$的所有奇点均位于复平面的左半平面$\text{Re}\ s<0$上。
\end{proposition}
\begin{proof}
    奇点在左半平面意味着单位冲激响应的增长速度至多相当于$\mathrm{e}^{-\alpha t}( \alpha>0)$，从而单位冲激响应绝对可积，系统稳定；反之，奇点在右半平面或虚轴上意味着单位冲激响应增长速度至少相当于$\mathrm{e}^{\alpha t}(\alpha\geq0)$，从而单位冲激响应不绝对可积，系统不稳定。
\end{proof}
\begin{definition}
    根据单位冲激响应的拉普拉斯变换的奇点位置，将因果线性时不变系统分为三类：
    \begin{itemize}
        \item \textbf{稳定系统}：所有奇点均位于左半平面上；
        \item \textbf{临界稳定系统}：至少有一个奇点位于虚轴上，且其余奇点均位于左半平面上；
        \item \textbf{不稳定系统}：至少有一个奇点位于右半平面上。
    \end{itemize}
\end{definition}

\subsection{初值定理与终值定理}\label{sub:initial_and_final_value}
本小节介绍\textbf{初值定理}和\textbf{终值定理}，它们可以通过拉普拉斯变换直接得到函数在$t=0$和$t\to\infty$时的极限值，而不需要求出函数的显式表达式。
\begin{theorem}[初值定理]
    设$f\in PS[0,\infty),f,f'\in\mathcal{E},F(s)=\mathcal{L} f(s)$，则
    \[f(0)=\lim_{\text{Re}\ s\to\infty}sF(s)\]
\end{theorem}
这表明系统的高频响应对应时域的瞬时的初始行为。
\begin{proof}
    由拉普拉斯变换的时域微分性质，有
    \[sF(s)=s\mathcal{L} f(s)=\mathcal{L} f'(s)+f(0)\]
    取极限$\text{Re}\ s\to\infty$，有
    \[\lim_{\text{Re}\ s\to\infty}sF(s)=\lim_{\text{Re}\ s\to\infty}\mathcal{L} f'(s)+f(0)=\lim_{\text{Re}\ s\to\infty}\int_{0}^{\infty}f'(t)\mathrm{e}^{-st}\,dt+f(0)\]
    由于$f'\in\mathcal{E} $，在$\text{Re}\ s$充分大时时，$f'(t)\mathrm{e}^{-st}\leq C \mathrm{e}^{-\epsilon t},\epsilon$为一充分小正数。由控制收敛定理，有
    \[\lim_{\text{Re}\ s\to\infty}\int_{0}^{\infty}f'(t)\mathrm{e}^{-st}\,dt=\int_{0}^{\infty}\lim_{\text{Re}\ s\to\infty}f'(t)\mathrm{e}^{-st}\,dt=0\]
    从而\[f(0)=\lim_{\text{Re}\ s\to\infty}sF(s)\]
\end{proof}

\begin{theorem}[终值定理]
    设$f\in\mathcal{E} $，$F(s)=\mathcal{L} f(s)$在$s=0$处有简单极点，且其余奇点均位于左半平面$\text{Re}\ s<0$上，则
    \[\lim_{t\to\infty}f(t)=\lim_{s\to 0}sF(s)\]
\end{theorem}
这表明系统的低频响应对应时域的稳态行为。
\begin{remark}
多数文献中，这个定理的允许$s=0$不是奇点，但由于$\sigma_s<0$时，多数情况下有$\sigma_c<0$，从而$f(t)\to 0$，因此这里我们只考虑非平凡情形。此外高阶极点对应的$\lim_{t\to\infty}f(t)=\infty$同样可以用以下方法证明。
\end{remark}
\begin{proof}
    考虑$F(s)$在$s=0$处的洛朗级数展开：
    \[F(s)=\sum_{n=-1}^{\infty}a_n s^n=\frac{a_{-1}}{s}+G(s),a_{-m}\neq 0\]
    其中$G(s)$在$s=0$的邻域内全纯，因此也在包含虚轴的某个右半平面$\text{Re}\ s>-\epsilon$上全纯，其拉普拉斯逆变换$g(t)$满足$g(t)\to 0,t\to\infty$，因此有\begin{align*}
        \lim_{s\to 0}sF(s)&=\lim_{s\to 0}\left(a_{-1}+sG(s)\right)=a_{-1}
    \end{align*}
    时域上，\begin{align*}
        f(t)&=\mathcal{L}^{-1} F(s)=\mathcal{L}^{-1}\left(\frac{a_{-1}}{s}\right)+\mathcal{L}^{-1} G(s)=a_{-1}u(t)+g(t)\to a_{-1},t\to\infty
    \end{align*}
    于是$\lim_{t\to\infty}f(t)=a_{-1}=\lim_{s\to 0}sF(s)$。
\end{proof}

%\subsection{在双边拉普拉斯变换中的应用}\label{sub:bilateral_laplace}
%本小节对\ref{sub:introduction_of_laplace}中提到的双边拉普拉斯变换的收敛域进行简单讨论，以便后面与z变换进行比较。双边拉普拉斯变换具有一些与单边拉普拉斯变换不同的性质（尤其是时域微分性质），但这不是本书的重点，读者可以仿照单边拉普拉斯变换自行推导它们。如有需要，读者可以暂时跳过本小节，在阅读\ref{sub:relation_z_laplace}时再回过头来阅读。

%回忆双边拉普拉斯变换的定义：
%\begin{definition}[双边拉普拉斯变换]
%    设函数$f$局部可积，$\exists a,b\in\mathbb{R},a<b$（我们允许$a=-\infty$和$b=\infty$），使得$f(t)\mathrm{e}^{-at}$和$f(t)\mathrm{e}^{-bt}$均绝对可积，则对$\text{Re}\ s\in[a,b]$，$f$的双边拉普拉斯变换定义为
%    \[\mathcal{L}f(s)=\int_{-\infty}^{\infty}f(t)\mathrm{e}^{-st}\,dt\]
%\end{definition}

%为了确定$a,b$的取值范围，我们可以拆分双边拉普拉斯积分，并仿照单边拉普拉斯变换的收敛域分析方法，分别讨论$t\to\infty$和$t\to-\infty$时$f(t)$的增长速率：
%\[\mathcal{L}f(s)=\int_{-\infty}^{0}f(t)\mathrm{e}^{-st}\,dt+\int_{0}^{\infty}f(t)\mathrm{e}^{-st}\,dt\]
%\begin{definition}[收敛横坐标]
%    \quad\\
%    \textbf{收敛横坐标}为使得拉普拉斯积分$\mathcal{L}f(s)=\int_{-\infty}^{\infty}f(t)\mathrm{e}^{-st}\,dt$最大收敛域的收敛轴实部（横坐标）$\sigma_{c}^-,\sigma_{c}^+$，使得$\mathcal{L}f(s)$恰好在$\text{Re}\ s\in(\sigma_{c}^-,\sigma_{c}^+)$上收敛，而在$\text{Re}\ s<\sigma_{c}^-$和$\text{Re}\ s>\sigma_{c}^+$上发散，即
%    \[\sigma_{c}^-=\sup\{a\in\mathbb{R}|\mathcal{L} _b f(s)\text{在}\text{Re}\ s\in(a,a+\epsilon)\text{上收敛}\}\]
%    \[\sigma_c^+=\inf\{b\in\mathbb{R}|\mathcal{L} _b f(s)\text{在}\text{Re}\ s\in(b-\epsilon,b)\text{上收敛}\}\]
%   这里要求$\text{Re}\ s\in(a,a+\epsilon)$和$\text{Re}\ s\in(b-\epsilon,b)$是为了处理收敛域为带状域的情况。
%\end{definition}
%要使$\int_{0}^{\infty}f(t)\mathrm{e}^{-st}\,dt$收敛，自然要求
%\[\text{Re}\ s>\varlimsup_{t\to\infty}\frac{ln|f(t)|}{t}\]
%因此可以看出$\sigma_c^-=\varlimsup_{t\to\infty}\frac{ln|f(t)|}{t}$。对于$\int_{-\infty}^{0}f(t)\mathrm{e}^{-st}\,dt$，记
%\[|f(t)|\sim C \mathrm{e}^{b t}\iff \lim_{t\to -\infty}\frac{|f(t)|}{C \mathrm{e}^{bt}}=1\]
%则有
%\[|f(t)|\sim C \mathrm{e}^{b t}\iff ln|f(t)|\sim b t + ln C\iff \frac{|f(t)|}{t}\sim b+\frac{ln C}{t}=b\]
%这时有
%$$\text{Re}\ s=b-\epsilon,\int_{-\infty}^{0}f(t)\mathrm{e}^{-st}\,dt\leq\int_{-\infty}^{0}C \mathrm{e}^{\epsilon t}\,dt<\infty$$
%因此可以看出$\sigma_c^+=\varliminf_{t\to -\infty}\dfrac{ln|f(t)|}{t}$（可以带入$f(t)=\mathrm{e}^{-t\sin t}$验证这里必须取下极限）。

%综上所述，我们得到：\begin{theorem}
%    设函数$f$局部可积，则其双边拉普拉斯变换的收敛横坐标可由以下公式得到：
%    \[\sigma_{c}^-=\varlimsup_{t\to\infty}\frac{ln|f(t)|}{t},\sigma_{c}^+=\varliminf_{t\to -\infty}\frac{ln|f(t)|}{t}\]
%\end{theorem}

%与单边拉普拉斯变换一样，我们允许$\sigma_{c}^-,\sigma_{c}^+=\pm\infty$。如果$\sigma_{c}^-=-\infty,\sigma_c^+=\infty$，则$\mathcal{L} _b f(s)$为整函数；如果$\sigma_{c}^-\geq\sigma_c^+$，则$\mathcal{L} _b f(s)$无定义。

%特别地，我们指出：\begin{itemize}
%    \item 如果$f(t)=0,t<0$，则$\sigma_{c}^-=-\infty$，双边拉普拉斯变换退化为单边拉普拉斯变换；
%    \item 如果$f(t)=0,t>0$，则$\sigma_{c}^+=\infty$；
%    \item 如果$f$紧支，则$\sigma_{c}^-=-\infty,\sigma_c^+=\infty$，双边拉普拉斯变换为整函数。
%    \item 如果$f(t)$在$t\to\pm\infty$时均与$\mathrm{e}^{a t}$表现类似，其双边拉普拉斯变换的收敛域为空集。除了前面三条中讨论的情况，典型的具有非空收敛域的例子如$f(t)=\mathrm{e}^{-|t|}$，其双边拉普拉斯变换的收敛域为$\text{Re}\ s\in(-1,1)$。
%\end{itemize}

%不难发现，双边拉普拉斯变换的收敛性分析与洛朗级数有着诸多相似之处，我们将在\ref{sub:relation_z_laplace}中讨论这一点。
